problem:  
Let \(M^n\) be a closed Riemannian manifold with nonnegative sectional curvature.  
By the Cheeger–Gromoll Structure Theorem, a finite cover of \(M\) is diffeomorphic to \(T^k \times N^{n-k}\), where \(T^k\) is a flat torus and \(N\) is compact, simply connected, and nonnegatively curved. The manifold \(M\) can be written as a quotient
\[
M = (T^k \times N)/\Gamma,
\]
where \(\Gamma\) is a finite group acting freely and isometrically, and the induced action on the flat torus \(T^k\) defines a holonomy representation  
\[
\rho: \Gamma \to \mathrm{GL}(k, \mathbb{Z}) \subset O(k).
\]  
The cohomology of \(M\) can be viewed as the \(\Gamma\)-invariant part of the cohomology of \(T^k \times N\),
\[
H^*(M) \cong H^*(T^k \times N)^\Gamma.
\]  

**Open problem (Curvature–Cohomology Rigidity Conjecture):**  
Among all possible finite free isometric \(\Gamma\)-actions arising from nonnegatively curved metrics, does there exist a universal constant \(c_n>0\), depending only on the dimension \(n\), such that  
\[
\frac{\dim H^*(M; \mathbb{Q})}{\dim H^*(T^k \times N; \mathbb{Q})} \ge c_n?
\]  
If not, determine the minimal possible ratio achievable under genuine geometric realizations of nonnegative sectional curvature. In particular, are there curvature-induced restrictions on the holonomy representation \(\rho\) that enforce a strictly positive lower bound for the invariant cohomology dimension?

Why is it a "good" problem:  
This problem explores an uncharted frontier at the intersection of curvature geometry, topology, and representation theory. Nonnegatively curved manifolds have a highly constrained metric and topological structure, yet the quantitative loss of cohomology under the transition from \(T^k \times N\) to a quotient \(M = (T^k \times N)/\Gamma\) remains fundamentally mysterious.  

Unlike purely algebraic settings, here the finite group \(\Gamma\) must arise from an isometric action compatible with nonnegative curvature. Understanding how curvature limits the possible representations \(\rho(\Gamma)\) is a deep geometric constraint problem that connects Riemannian geometry and equivariant cohomology.  

A positive resolution would uncover a new rigidity phenomenon—showing that curvature forbids arbitrarily “twisted” cohomological collapses. Even if the inequality fails in general, identifying geometric obstructions or constructing counterexamples would substantially advance our understanding of how local curvature constraints influence global cohomological invariants.  

The statement is concise and concrete, yet its resolution requires integrating differential geometry, algebraic topology, and representation theory—an exemplary “narrow entrance, profound depth” research problem.